\section{Low-rank preserves RKHS}\label{sec:no-rkhs-advantage}
In the NTK regime, training dynamics linearize: \( f_t(x) \approx f_0(x) + \langle \nabla_\theta f_0(x), \theta_t - \theta_0 \rangle \). The learned function stays in the RKHS induced by \( \Theta \), with norm:
\[
\|f\|_{\mathcal{H}_\Theta}^2 = \langle f, \Theta^{-1} f \rangle_{L^2}.
\]
For RF-LR with frozen random features, the RKHS is characterized as follows. Universal approximation holds due to full support of \( w^{(\ell)} \) (frozen Gaussian draws), so any function in the Barron space~\cite{e2021barronspaceflowinducedfunction} can be approximated with error \( O(1/\sqrt{N}) \). 
So the RKHS is the same as for MLPs NTK analysis shows weights are nearly frozen through training. We tackle directly this remark by freezing the weights instead of performing weights-deviation analysis on MLPs.
This also removes the EOC condition on the bias.

We can now state a result highly similar to ~\cite{bietti2021deepequalsshallowrelu}: even for in the extensive rank regime, RF-LR NTK under the NTK perspective, deeper architectures do not yield a larger RKHS than shallow networks.

\vspace{0.5em}
\subsection{Rank-driven concentration in RKHS (Gaussian norms product)}\label{subsec:rank-driven-concentration}
Having established that low rank and depth confers no RKHS advantage, we now turn to quantifying how the low-rank structure controls kernel fluctuations. We can show that the radial component of the RF-LR NTK concentrates exponentially fast in the rank \( r \), yielding high-dimensional control of RKHS fluctuations. To state this result, we first require a homogeneity assumption on the activation function.

\begin{assumption}[Homogeneous activation]
The activation function \( \sigma \) is positively homogeneous: \( \sigma(\alpha z) = \alpha \sigma(z) \) for all \( \alpha \ge 0 \) and \( z \in \mathbb{R} \). In particular, ReLU satisfies this property.
\end{assumption}

We can now state our main concentration result.

\begin{theorem}[Concentration in rank for homogeneous activation]\label{thm:rkhs_rank_concentration}
Under homogeneous activation assumption (Assumption~\ref{assump:homogeneous_activation}), let \( x_1, y_1 \in \mathbb{R}^r \) be the rank-\( r \) Gaussian projections of inputs \( x, y \), with arbitrary fixed correlation \( \rho \in [-1,1] \). Define
\[
W \;=\; \frac{\|x_1\|\,\|y_1\|}{r}.
\]
Then for any \( \epsilon \in (0,1) \),
\[
\mathbb{P}\!\left(\left|W-1\right|\;\ge\;\epsilon\right)\;\le\;4\,\exp\!\left(-\frac{r\,\epsilon^2}{8}\right).
\]
Moreover, for larger deviations \( \epsilon \ge 1 \), there exist absolute constants \( c_1,c_2>0 \) such that
\[
\mathbb{P}\!\left(\left|W-1\right|\;\ge\;\epsilon\right)\;\le\;c_1\,\exp\!\left(-c_2\,r\,\epsilon\right).
\]
In particular, the radial component of the RF-LR NTK concentrates exponentially fast in \( r \), yielding high-dimensional control of RKHS fluctuations.
\end{theorem}

We now provide an intuitive explanation of the proof strategy.

\paragraph{Proof intuition.}
Set \( U=\|x_1\|/\sqrt{r} \) and \( V=\|y_1\|/\sqrt{r} \). \( U \) and \( V \) concentrate around 1 with sub-Gaussian tails \( \exp(-r\delta^2) \). Using \( |UV-1|\le |U-1|+|V-1|+|U-1||V-1| \) and a union bound with \( \delta=\epsilon/2 \) yields the stated \( 4\exp(-r\epsilon^2/8) \) bound. The full proof is given in Appendix~\ref{appendix:proof-rkhs-rank-concentration}.
\begin{remark}
The result avoids direct manipulation of Kibble's bivariate chi-square density, relying instead on Gaussian isoperimetry for norms and a simple algebraic bound for products.
\end{remark}

\vspace{0.5em}
\subsection{three-layer NTK concentration}\label{subsec:three-layer-ntk-concentration}
We can now apply this concentration result to obtain bounds for the three-layer NTK.

\begin{corollary}[Concentration bound for three-layer NTK]\label{cor:two_layer_ntk_concentration}
Under Assumption~\ref{assump:homogeneous_activation}, let \( \hat{\Theta}^{(2)}_r(x,x') = \Psi(\hat{\rho}_r, \hat{w}_r) \) be the empirical three-layer NTK, where \( \hat{\rho}_r \) is the sample correlation and \( \hat{w}_r = \|x_1\|\|y_1\|/r \), with \( x_1, y_1 \in \mathbb{R}^r \) being the rank-\( r \) Gaussian projections of inputs \( x, y \) with population correlation \( \rho \in [-1,1] \). Let \( K_{\infty}(\rho) = \Theta^{(1)}(\rho)(1 - \arccos(\rho)/\pi) + \Sigma^{(1)}(\rho) + 1 \) be the deterministic kernel limit.

Then for any \( \epsilon \in (0,1) \), there exist constants \( C_1, C_2 > 0 \) depending on \( \rho \) and the kernel smoothness such that
\[
\mathbb{P}\!\left(\left|\hat{\Theta}^{(2)}_r(x,x') - K_{\infty}(\rho)\right|\;\ge\;\epsilon\right)\;\le\;C_1\,\exp\!\left(-C_2\,r\,\epsilon^2\right).
\]
\end{corollary}

Proof: See Appendix~\ref{appendix:proof-three-layer-ntk-concentration}.

\begin{remark}[Critical cancellation of \( 1/r \) corrections]
While the general concentration rate is \( O_p(1/\sqrt{r}) \) via Fisher's CLT~\cite{fisher1915probable,hotelling1953new}, the randomness induced from the gaussian process low rank output remove the \( 1/r \) correction term in \( \Sigma^{(2)} \) that would normally appear through the radial component \( \hat{w}_r \). This cancellation is critical: it is precisely the path where finite-width corrections \( 1/N \) analysis~\cite{dyer2019asymptoticswidenetworksfeynman} and Feynman diagram decay come into play, proving that the fluctuation rate is \( 1/r \) rather than \( 1/\sqrt{r} \).
\end{remark}

\begin{remark}[Extension to general \( L \) layers]
The extension of this concentration bound to general depth \( L \) layers under Assumption~\ref{assump:homogeneous_activation} requires controlling nested conditional means through \( L \) layers and managing a quadratic number \( O(L^2) \) of cross-terms arising from the recursive NTK structure. The technical challenge lies in tracking how the random correlations \( \rho_\ell \) and norm products \( \|h^{(\ell)}(x)\|\|h^{(\ell)}(x')\| \) propagate through layers while maintaining exponential concentration rates. This analysis is deferred to a future version.
\end{remark}

\vspace{0.5em}
\subsection{Mean NTK over Fisher and Kibble distributions}\label{subsec:conditional-mean-ntk}
Having established concentration results, we now turn to the mean NTK. The key observation is that the empirical NTK depends on random variables that follow Fisher~\cite{fisher1915probable,hotelling1953new} and Kibble distributions. The mean NTK is obtained by taking expectations over these distributions.

For the three-layer case, the empirical NTK is:
\[
\hat{\Theta}^{(2)}_r(x,x') = \Theta^{(1)}(\rho) \cdot \left(1 - \frac{\arccos(\hat{\rho}_r)}{\pi}\right) + \frac{1}{r} \cdot \Sigma^{(1)}(\hat{\rho}_r) \cdot \|x_1\|\|y_1\| + 1,
\]
where \( \hat{\rho}_r \) follows Fisher's distribution~\cite{fisher1915probable,hotelling1953new} (centered at \( \rho \)) and \( \|x_1\|\|y_1\| \) follows Kibble's distribution. The mean NTK is therefore:
\[
\tilde{\Kop}^{(2)}(\rho) = \mathbb{E}_{\text{Fisher}, \text{Kibble}}\left[\hat{\Theta}^{(2)}_r(x,x')\right] = \Theta^{(1)}(\rho) \cdot \mathbb{E}\!\left[1 - \frac{\arccos(\hat{\rho}_r)}{\pi}\right] + \frac{1}{r} \cdot \mathbb{E}\!\left[\Sigma^{(1)}(\hat{\rho}_r) \cdot \|x_1\|\|y_1\|\right] + 1,
\]
where the expectations are taken over the Fisher distribution of \( \hat{\rho}_r \) and the Kibble distribution of \( (\|x_1\|^2, \|y_1\|^2) \).

\begin{corollary}[Mean NTK over Fisher–Kibble has same RKHS]\label{cor:conditional_rkhs_homogeneous}
Under Assumption~\ref{assump:homogeneous_activation} and the RF-LR setting with ReLU nonlinearity and isotropic random features on \( \mathbb{S}^{d-1} \), the mean NTK \( \tilde{\Kop}^{(L)} \) (obtained by taking expectations over Fisher and Kibble distributions at each layer) is a zonal kernel on \( \mathbb{S}^{d-1} \) that induces the same RKHS as the shallow (three-layer) ReLU kernel. In particular, the mean NTK RKHS coincides with the unconditional RKHS as sets with equivalent norms.

For the three-layer case (\( L=2 \)), the mean NTK is:
\[
\tilde{\Kop}^{(2)}(\rho) = \Theta^{(1)}(\rho) \cdot \mathbb{E}_{\hat{\rho}_r \sim \text{Fisher}(\rho,r)}\!\left[1 - \frac{\arccos(\hat{\rho}_r)}{\pi}\right] + \frac{1}{r} \cdot \mathbb{E}_{\hat{\rho}_r, \|x_1\|, \|y_1\|}\!\left[\Sigma^{(1)}(\hat{\rho}_r) \cdot \|x_1\|\|y_1\|\right] + 1,
\]
where the expectations are over Fisher~\cite{fisher1915probable,hotelling1953new} and Kibble distributions. This yields a zonal kernel with a Puiseux expansion that differs from the deterministic kernel \( K_{\infty}(\rho) \), but preserves the same endpoint behavior, ensuring RKHS equivalence.
\end{corollary}

\paragraph{Intuition and technical challenges.}
This corollary is highly non-trivial because it requires computing expectations of \( \arccos(\hat{\rho}_r) \) where \( \hat{\rho}_r \sim \text{Fisher}(\rho,r) \)~\cite{fisher1915probable,hotelling1953new}, which involves integrating over Fisher's full distribution density near the boundary \( \rho = 1 \). The key technical challenge is analyzing the Puiseux expansion of the mean NTK near \( \rho = 1 \), where both \( \arccos(\hat{\rho}_r) \) and Fisher's density exhibit singular behavior.

The proof (Appendix~\ref{appendix:proof-conditional-rkhs-homogeneous}) computes the integral:
\[
\mathbb{E}[\arccos(\hat{\rho}_r)] \;=\; \int_{-1}^{1} \arccos(u)\;
\frac{(r-2)\,\Gamma(r-1)\,\big(1-(1-t)^2\big)^{\frac{r-1}{2}}\,(1-u^2)^{\frac{r-4}{2}}}{\sqrt{2\pi}\,\Gamma\!\left(r-\tfrac{1}{2}\right)\,\big(1-(1-t)\,u\big)^{\,r-\frac{3}{2}}}\;
{}_2F_1\!\left(\tfrac{1}{2},\,\tfrac{1}{2};\, r-\tfrac{1}{2};\, \tfrac{1+(1-t)\,u}{2}\right)\; du,
\]
Near this boundary, we rely on asymptotic expansions of hypergeometric functions near their argument \( z \to 1^- \) (see~\cite{nist_dlmf} and~\cite{abramowitz1964handbook}).

The crucial insight is that while the mean NTK has a different Puiseux expansion than the deterministic kernel \( K_{\infty}(\rho) \), the leading-order \( t^{1/2} \) term (which determines spectral decay and RKHS structure) is preserved. This preservation of endpoint behavior under Fisher–Kibble expectations is what ensures RKHS equivalence, despite the different higher-order corrections that appear in the mean kernel.


\vspace{0.5em}
\subsection{Spectral decay and Puiseux expansions}\label{subsec:spectral-decay-puiseux}
We now apply the spectral decay theory of~\cite{bietti2021deepequalsshallowrelu} to characterize the RKHS of the RF-LR NTK via endpoint expansions. The base kernel \( \Sigma^{(1)} \) and derivative kernel \( \dot{\Sigma}^{(1)} \) are the fundamental building blocks for the NTK recursion.

From the proof, near \( \rho = 1 \) (for \( t \ge 0 \)), the mean two-layer NTK admits the Puiseux expansion (see Appendix~\ref{appendix:proof-conditional-rkhs-homogeneous}):
\[
\tilde{\Kop}^{(2)}(1-t)
\;=\;
K_\infty(1)
\;-\;
\Bigg[\frac{2}{\pi} \;+\; \frac{\sqrt{2}}{2\pi}\,I(r)\Bigg]\,t^{1/2}
\;+\; O(t^{3/2}),
\]
where the leading non-polynomial term has exponent \( \nu = 1/2 \) with coefficient \( c_1(r) = -\big[2/\pi + (\sqrt{2}/(2\pi))\,I(r)\big] \). Similarly, near \( \rho = -1 \), by rotational symmetry and homogeneity of ReLU, the expansion has the same structure with \( c_{-1}(r) = c_1(r) \) (see~\cite{bietti2021deepequalsshallowrelu}).

We can now apply Theorem 1 of~\cite{bietti2021deepequalsshallowrelu} to obtain the spectral decay characterization:

\begin{theorem}[Spectral decay for mean RF-LR NTK with final constant]\label{cor:spectral_decay_bietti}
Let \( \kappa(\rho) = \tilde{\Kop}^{(2)}(\rho) \) be the mean NTK on \( \mathbb{S}^{d-1} \). From the Puiseux expansion in Appendix (see the explicit constant at the end of Appendix~\ref{appendix:proof-conditional-rkhs-homogeneous}), \( \kappa \) admits, for \( t\ge 0 \),
\[
\kappa(1-t) = p_1(t) + c_1(r)\, t^{1/2} + o(t^{1/2}),\qquad
\kappa(-1+t) = p_{-1}(t) + c_{-1}(r)\, t^{1/2} + o(t^{1/2}),
\]
with exponent \( \nu = 1/2 \). For ReLU and normalized inputs,
\[
c_1(r) \;=\; -\left[\frac{2}{\pi} \;+\; \frac{\sqrt{2}}{2\pi}\,I(r)\right],
\qquad
I(r) \;=\; \frac{(r-2)\,2^{\,r-\frac{5}{2}}\,\Gamma\!\left(\frac{r-1}{2}\right)\Gamma\!\left(\frac{r}{2}-1\right)}{\sqrt{2\pi}\,\Gamma(r-1)}\cdot\frac{\Gamma\!\left(r-\tfrac{1}{2}\right)\Gamma\!\left(r-\tfrac{3}{2}\right)}{\Gamma(r-1)^2},
\]
and, by symmetry of the ReLU kernel on \( \mathbb{S}^{d-1} \), \( c_{-1}(r) = c_1(r) \).

Then by Theorem 1 of~\cite{bietti2021deepequalsshallowrelu}, the eigenvalues \( \mu_k \) (in the spherical harmonics basis) satisfy:
\[
\mu_k \sim C(d,\nu)\,k^{-d-2\nu+1} = C(d)\,k^{-d-2(1/2)+1} = C(d)\,k^{-d},
\]
where \( C(d) > 0 \) depends only on \( d \). Since \( c_1(r) = c_{-1}(r) \), we have \( |c_1(r)| = |c_{-1}(r)| \), hence:
\begin{itemize}
    \item For even \( k \): \( \mu_k \sim \big(c_1(r)+c_{-1}(r)\big)\,C(d)\,k^{-d} = 2\,c_1(r)\,C(d)\,k^{-d} \).
    \item For odd \( k \): since \( c_1(r) - c_{-1}(r) = 0 \), \( \mu_k = o(k^{-d}) \) (faster-than-polynomial in this parity).
\end{itemize}
Therefore, the dominant eigenvalue decay rate is \( k^{-d} \), determined by the even-parity eigenvalues.
\end{theorem}

\begin{corollary}[No advantage in RKHS]\label{thm:no_rkhs_advantage}
Under the RF-LR setting with ReLU nonlinearity, isotropic random features on the sphere \( \mathbb{S}^{d-1} \), and in the kernel (NTK) regime, the depth-\(L\) NTK induces the \emph{same} RKHS as the shallow (three-layer) ReLU kernel. In particular, the corresponding RKHSs coincide as sets with equivalent norms. This is a direct application of Theorem 1 in~\cite{bietti2021deepequalsshallowrelu}.
\end{corollary}

\begin{remark}
This theorem and corollary are direct applications of~\cite{bietti2021deepequalsshallowrelu} to the RF-LR setting. Theorem~\ref{cor:spectral_decay_bietti} supplies explicit rates for approximation and generalization in the RF-LR NTK RKHS and, together with the explicit kernel formulas above, justifies that deeper NTKs do not enlarge the RKHS beyond the shallow ReLU case (Corollary~\ref{thm:no_rkhs_advantage}).
\end{remark}
